\documentclass[11pt,a4paper]{article}

% --- Core Packages ---
\usepackage[utf8]{inputenc}
\usepackage[T1]{fontenc}
\usepackage{lmodern}
\usepackage[margin=1in]{geometry}
\usepackage{xcolor}
\usepackage{graphicx}
\usepackage{booktabs}
\usepackage{amsmath,amssymb}
\usepackage{listings}
\usepackage{fancyhdr}
\usepackage{enumitem}
\usepackage{tabularx}
\usepackage{microtype}
\usepackage[hidelinks,colorlinks=true,linkcolor=blue!70!black,urlcolor=blue!80!black,citecolor=blue!70!black]{hyperref}

% --- Color Palette ---
\definecolor{codegreen}{rgb}{0,0.6,0}
\definecolor{codegray}{rgb}{0.5,0.5,0.5}
\definecolor{codepurple}{rgb}{0.58,0,0.82}
\definecolor{backcolour}{rgb}{0.96,0.97,0.98}
\definecolor{primary}{RGB}{18,52,86}
\definecolor{alertred}{RGB}{180,30,30}

% --- Code Listings Configuration ---
\lstdefinestyle{mystyle}{
    backgroundcolor=\color{backcolour},
    commentstyle=\color{codegreen},
    keywordstyle=\color{primary}\bfseries,
    numberstyle=\tiny\color{codegray},
    stringstyle=\color{codepurple},
    basicstyle=\ttfamily\footnotesize,
    breakatwhitespace=false,
    breaklines=true,
    captionpos=b,
    keepspaces=true,
    numbers=left,
    numbersep=7pt,
    showspaces=false,
    showstringspaces=false,
    showtabs=false,
    tabsize=2,
    frame=single,
    rulecolor=\color{black!20}
}
\lstset{style=mystyle}

% --- Header and Footer ---
\setlength{\headheight}{14pt}
\pagestyle{fancy}
\fancyhf{}
\fancyhead[L]{\textsc{IoT Wi-Fi Network Security \& SOC Deployment}}
\fancyhead[R]{\thepage}
\fancyfoot[C]{\footnotesize AI-Assisted Real-Time Intrusion Detection System $\cdot$ Operational Manual}
\renewcommand{\headrulewidth}{0.5pt}
\renewcommand{\footrulewidth}{0.4pt}

% --- Document Metadata ---
\title{\vspace{-1.5cm}\textbf{\LARGE AI-Assisted IoT Wi-Fi Network Security and Anomaly Detection System}\\[0.4cm]
\Large Complete Operational Manual, Firmware Flashing Guide, and SOC Frontend Deployment}
\author{\textbf{Justin \& Research Team} \\ Department of Computing and Network Security}
\date{November 2026}

\begin{document}

\maketitle
\thispagestyle{empty}

\begin{abstract}
This operational manual provides exhaustive, step-by-step instructions for deploying, configuring, and operating the AI-Assisted IoT Wi-Fi Network Security and Anomaly Detection System. It covers: (1) compiling and flashing the dual-sensor and promiscuous management sniffer firmware onto the ESP32 microcontroller using the Arduino IDE; (2) required hardware pinouts, power distribution, and sensor calibration; (3) one-click batch execution of the Streamlit Security Operations Center (SOC) frontend via \texttt{run\_frontend.bat}; (4) running the background ingestion daemon (\texttt{live\_collector.py}); and (5) diagnostic workflows for hardware channel auto-alignment, memory safety, and attack alert investigation.
\end{abstract}

\vspace{0.3cm}
\tableofcontents
\vspace{0.5cm}
\hrule
\vspace{0.5cm}

% =============================================================================
\section{System Architecture Overview}
% =============================================================================

The system is structured as an out-of-band, multi-stage hybrid detection platform. It captures both digital packet telemetry over the local Wi-Fi local area network (via TShark) and raw 802.11 management attack frames floating through physical air space (via ESP32 promiscuous sniffing), correlated with environmental IoT sensors.

\begin{figure}[htbp]
\centering
\begin{verbatim}
       +--------------------------------------------------------+
       |             2.4 GHz Physical Radio Airspace            |
       +--------------------------------------------------------+
              | (802.11 Deauth/Disassoc)       | (Wi-Fi 802.11 b/g/n)
              v                                v
       +--------------+                 +--------------------+
       | ESP32 Node   |                 | Windows Hotspot    |
       | (Sniffer)    |                 | Gateway Adapter    |
       +--------------+                 +--------------------+
              |                                   |
              | USB Serial UART (115200 baud)     | TShark Packet Stream
              v                                   v
       +--------------------------------------------------------+
       |               Python Background Ingestion               |
       |                   (src/live_collector.py)               |
       |                                                        |
       |  - 1-Second Sliding Window Feature Extraction          |
       |  - Frozen Dual-Detector Isolation Forest Inference      |
       |  - Deterministic Network Attack Rules (DoS, Scan, ARP)  |
       |  - Dynamic ESP32 Sniffer Channel Auto-Alignment         |
       +--------------------------------------------------------+
                                   |
                                   | WAL Journal Writes
                                   v
       +--------------------------------------------------------+
       |            SQLite Database (data/live_monitoring.db)    |
       +--------------------------------------------------------+
                                   |
                                   | 2-Second Polling Fragment
                                   v
       +--------------------------------------------------------+
       |        Streamlit SOC Frontend (src/live_dashboard.py)   |
       |                 Launched via: run_frontend.bat         |
       +--------------------------------------------------------+
\end{verbatim}
\caption{End-to-End Hybrid Detection and Telemetry Ingestion Architecture.}
\end{figure}

\subsection{Key Design Guarantee: Zero Network Footprint}
A fundamental architectural principle of this system is that \textbf{the ESP32 microcontroller does not join any Wi-Fi network and does not host an Access Point (SoftAP)}. 
\begin{itemize}[noitemsep]
    \item It has \textbf{no IP address}, no MAC binding on the host router, and zero open TCP/UDP sockets.
    \item It listens passively to raw 802.11 radio frames on its configured channel using ESP-IDF promiscuous filters.
    \item All telemetry is delivered out-of-band via a physical \textbf{USB Serial cable} directly into the host PC, ensuring that wireless jamming or deauthentication attacks cannot silence or compromise sensor telemetry.
\end{itemize}

% =============================================================================
\section{Hardware Requirements and Sensor Wiring}
% =============================================================================

\subsection{Bill of Materials}
\begin{enumerate}[noitemsep]
    \item \textbf{ESP32 Development Board}: ESP32-WROOM-32, NodeMCU-32S, or ESP32 DevKit v1.
    \item \textbf{DHT11 Sensor}: Digital Temperature and Relative Humidity sensor (3-pin module or 4-pin bare sensor).
    \item \textbf{AM312 PIR Sensor}: Miniature Passive Infrared Human Motion detector.
    \item \textbf{Host Computer}: Windows 10/11 with Wi-Fi adapter supporting Windows Mobile Hotspot.
    \item \textbf{Micro-USB / Type-C Data Cable}: High-quality cable carrying power and D+/D- data lines.
\end{enumerate}

\subsection{Pinout and Wiring Table}
All connections use standard 3.3V logic levels. \textbf{Do NOT connect sensors to the 5V / VIN pin}, as overvoltage will damage the ESP32 GPIO inputs.

\begin{table}[htbp]
\centering
\begin{tabularx}{\textwidth}{llllX}
\toprule
\textbf{Component} & \textbf{Pin Label} & \textbf{ESP32 Pin} & \textbf{Voltage} & \textbf{Notes / Function} \\
\midrule
\textbf{DHT11} & VCC (+) & \textbf{3V3} & 3.3V & Power supply rail \\
\textbf{DHT11} & GND ($-$) & \textbf{GND} & 0V & Common ground reference \\
\textbf{DHT11} & DATA / OUT & \textbf{GPIO 4} & 3.3V Logic & Single-wire digital bidirectional data \\
\midrule
\textbf{AM312 PIR} & VCC & \textbf{3V3} & 3.3V & Power supply rail (3.0V -- 5.0V compatible) \\
\textbf{AM312 PIR} & GND & \textbf{GND} & 0V & Common ground reference \\
\textbf{AM312 PIR} & OUT & \textbf{GPIO 27} & 3.3V Logic & Digital high (1) on motion, low (0) on idle \\
\bottomrule
\end{tabularx}
\caption{Complete ESP32 Sensor Hardware Pin Interconnections.}
\end{table}

\paragraph{Pull-Up Resistor Notice:} If you are using a bare 4-pin DHT11 sensor rather than a 3-pin breakout module, place a $4.7\text{ k}\Omega$ to $10\text{ k}\Omega$ resistor between DHT11 pin 2 (DATA) and pin 1 (3V3). Breakout boards with 3 pins already include this SMD resistor on the PCB.

% =============================================================================
\section{ESP32 Firmware Setup and Flashing via Arduino IDE}
% =============================================================================

The firmware source code is located at:
\begin{center}
\texttt{firmware/esp32\_sensor\_node/esp32\_sensor\_node.ino}
\end{center}

Follow these exact steps to compile and flash the firmware onto the ESP32.

\subsection{Step 1: Install Arduino IDE}
Download and install \textbf{Arduino IDE 2.3.x} (or 1.8.19) from \url{https://www.arduino.cc/en/software}.

\subsection{Step 2: Install the ESP32 Board Package}
\begin{enumerate}
    \item Open the Arduino IDE.
    \item Navigate to \textbf{File $\rightarrow$ Preferences}.
    \item Locate the field labeled \textbf{Additional Boards Manager URLs} and paste the following official Espressif URL:
\begin{lstlisting}[numbers=none]
https://raw.githubusercontent.com/espressif/arduino-esp32/gh-pages/package_esp32_index.json
\end{lstlisting}
    \item Click \textbf{OK}.
    \item Open \textbf{Tools $\rightarrow$ Board $\rightarrow$ Boards Manager...} (or click the board icon on the left sidebar in Arduino IDE 2.x).
    \item Type \texttt{esp32} in the search bar.
    \item Locate \textbf{esp32 by Espressif Systems} and click \textbf{Install} (version 2.0.14 or later recommended).
\end{enumerate}

\subsection{Step 3: Install Required Sensor Libraries}
Open \textbf{Tools $\rightarrow$ Manage Libraries...} (or click the library book icon on the left sidebar).
\begin{enumerate}
    \item In the search box, type: \texttt{DHT sensor library}.
    \item Locate \textbf{DHT sensor library by Adafruit} and click \textbf{Install}.
    \item If prompted with ``\textit{Install All Dependencies?}'', click \textbf{Install All}. This automatically installs \textbf{Adafruit Unified Sensor}.
\end{enumerate}

\subsection{Step 4: Configure Board Settings in Arduino IDE}
Connect the ESP32 to your PC using the USB cable. Under the \textbf{Tools} menu, set:
\begin{itemize}[noitemsep]
    \item \textbf{Board:} \texttt{ESP32 Dev Module} (or \texttt{NodeMCU-32S})
    \item \textbf{Upload Speed:} \texttt{921600} (use \texttt{115200} if upload errors occur on budget USB cables)
    \item \textbf{CPU Frequency:} \texttt{240MHz (WiFi/BT)}
    \item \textbf{Flash Frequency:} \texttt{80MHz}
    \item \textbf{Flash Mode:} \texttt{QIO}
    \item \textbf{Port:} Select the detected COM port (e.g., \texttt{COM3}, \texttt{COM6}).
\end{itemize}

\subsection{Step 5: Compile and Upload Firmware}
\begin{enumerate}
    \item Open \texttt{firmware/esp32\_sensor\_node/esp32\_sensor\_node.ino} in the Arduino IDE.
    \item Click the \textbf{Upload} arrow icon ($\rightarrow$) or press \texttt{Ctrl + U}.
    \item \textbf{Bootloader Synchronization}: If the console shows:
\begin{lstlisting}[numbers=none]
Connecting........_____....._____.....
\end{lstlisting}
    immediately press and hold the physical \textbf{BOOT} (or \textbf{IO0}) button on the ESP32 board for 2 seconds until the upload progress percentage ($1\% \dots 100\%$) appears, then release it.
    \item The upload finishes with \texttt{Hash of data verified} and \texttt{Hard resetting via RTS pin...}.
\end{enumerate}

\subsection{Step 6: Hardware Verification via Serial Monitor}
\begin{enumerate}
    \item In the Arduino IDE, open \textbf{Tools $\rightarrow$ Serial Monitor} (\texttt{Ctrl + Shift + M}).
    \item Set the baud rate dropdown to \textbf{115200 baud} and line ending to \textbf{Newline}.
    \item Press the \textbf{EN / RST} button on the ESP32. You will observe the boot announcement followed by periodic telemetry:
\begin{lstlisting}[numbers=none]
{"status":"ESP32_TELEMETRY_INITIALIZED","baud":115200}
{"node":"esp32_gateway","seq":1,"uptime_ms":2014,"dht11":{"valid":1,"temp_c":24.8,"humidity":52.1},"pir":{"motion":0}}
\end{lstlisting}
    \item \textbf{CRITICAL}: \textbf{Close the Arduino Serial Monitor window} before starting the Python backend collector. Windows does not allow multiple processes to share the same COM port simultaneously.
\end{enumerate}

% =============================================================================
\section{One-Click Frontend Deployment (\texttt{run\_frontend.bat})}
% =============================================================================

The project includes an automated, self-healing batch launcher: \texttt{run\_frontend.bat}.

\subsection{How It Works Under the Hood}
When executed, \texttt{run\_frontend.bat} performs the following sequential checks:
\begin{enumerate}[noitemsep]
    \item \textbf{Project Root Anchoring}: Uses \texttt{\%~dp0} to guarantee correct working directory context regardless of whether it is double-clicked from Windows File Explorer or invoked via terminal.
    \item \textbf{Virtual Environment Detection}: Searches for \texttt{.venv\textbackslash Scripts\textbackslash python.exe}. If absent, it smoothly falls back to the system Python installation.
    \item \textbf{Autonomous Database Guard}: Checks if \texttt{data\textbackslash live\_monitoring.db} exists. If not found, it automatically calls \texttt{init\_live\_database(..., seed\_defaults=True)} to provision all 9 database tables and default MAC device profiles.
    \item \textbf{Streamlit SOC Launch}: Executes \texttt{streamlit run src\textbackslash live\_dashboard.py -{}-server.port 8501}.
    \item \textbf{Failure Trapping}: If any step exits abnormally, the window remains open with an explanatory error code rather than abruptly closing.
\end{enumerate}

\subsection{Execution Methods}

\paragraph{Method A: Double-Click (Recommended for Operators)}
Navigate to the project folder in Windows File Explorer and simply double-click:
\begin{center}
\fbox{\textbf{\texttt{run\_frontend.bat}}}
\end{center}

\paragraph{Method B: Command Prompt / PowerShell}
Open a terminal in the project directory and run:
\begin{lstlisting}[language=bash, numbers=none]
.\run_frontend.bat
\end{lstlisting}

\subsection{Accessing the Live SOC Dashboard}
Once the terminal displays:
\begin{lstlisting}[numbers=none]
  You can now view your Streamlit app in your browser.
  Local URL: http://localhost:8501
  Network URL: http://192.168.1.2:8501
\end{lstlisting}
open your web browser and navigate to:
\begin{center}
\textbf{\url{http://localhost:8501}}
\end{center}

% =============================================================================
\section{Running the Complete Production Pipeline}
% =============================================================================

To run the complete system with live hardware and real-time network traffic analysis, follow this standard operational sequence:

\subsection{Step 1: Enable Windows Mobile Hotspot}
\begin{enumerate}[noitemsep]
    \item Open Windows Settings $\rightarrow$ \textbf{Network \& Internet} $\rightarrow$ \textbf{Mobile hotspot}.
    \item Set Network band to \textbf{2.4 GHz} (802.11 b/g/n).
    \item Toggle Mobile Hotspot to \textbf{On}.
    \item Connect target IoT devices (e.g., test smartphones or ESP8266/ESP32 clients) to this hotspot.
\end{enumerate}

\subsection{Step 2: Start the Live Ingestion Collector}
Open a terminal window and run:
\begin{lstlisting}[language=bash, numbers=none]
.venv\Scripts\python.exe src/live_collector.py --port COM6
\end{lstlisting}
\textit{(Replace \texttt{COM6} with the COM port assigned to your ESP32 in Device Manager).}

The collector will:
\begin{itemize}[noitemsep]
    \item Query Windows for the active Mobile Hotspot virtual adapter and Wi-Fi channel.
    \item Connect to the ESP32, issue a cryptographic probe handshake (\texttt{\{"cmd":"probe"...\}}), and synchronize the ESP32 sniffer radio to the hotspot's channel.
    \item Launch TShark to capture 1-second packet windows.
    \item Re-create the frozen Session 1 Dual-Detector bundle and begin real-time scoring.
\end{itemize}

\subsection{Step 3: Launch the Frontend}
In a separate terminal (or by double-clicking), execute:
\begin{lstlisting}[language=bash, numbers=none]
.\run_frontend.bat
\end{lstlisting}
The dashboard will display incoming sensor telemetry, network traffic rates, and attack alerts with 2-second streaming auto-refresh.

% =============================================================================
\section{Dashboard Features and Alert Interpretation}
% =============================================================================

\subsection{SOC Telemetry Panels}
\begin{itemize}
    \item \textbf{Collector Health Matrix}: Displays real-time heartbeats, database WAL file size, serial connection backoff state, and hardware channel alignment.
    \item \textbf{Dual-Detector Engine}: Renders 1-second Isolation Forest outlier scores ($\pm 1$) and packet timing distributions for each enrolled device.
    \item \textbf{IoT Environmental Correlation}: Charts live temperature (°C), humidity (\%), and PIR human motion states alongside network packet surges.
\end{itemize}

\subsection{Attack Classification Matrix}
\begin{table}[htbp]
\centering
\begin{tabularx}{\textwidth}{p{3.2cm}p{5.0cm}X}
\toprule
\textbf{Attack Type} & \textbf{Detection Logic} & \textbf{Primary Indicators} \\
\midrule
\textbf{DoS / Flood} & Sustained high packet-per-second rate targeted at a specific internal node or gateway. & Rates $\ge 60\text{ PPS}$ over $\ge 2\text{ s}$, high concentration index ($>0.70$). Suppressed if entire fleet surges simultaneously. \\
\addlinespace
\textbf{Horizontal Scan} & Single scanner probing identical ports across multiple LAN IP targets. & $\ge 5$ distinct destination IPs targeted with SYN/UDP frames within sliding window. \\
\addlinespace
\textbf{Vertical Scan} & Single scanner probing multiple destination ports on a single IP target. & $\ge 8$ distinct transport ports probed on target host within sliding window. \\
\addlinespace
\textbf{IP / MAC Spoofing} & Same IP claimed by an unregistered or conflicting MAC address. & Conflicting IP-to-MAC ARP bindings against the enrolled device registry. \\
\addlinespace
\textbf{Wi-Fi Deauthentication} & Over-the-air 802.11 management frames intercepted by ESP32 sniffer. & Subtype 12 (Deauth) or Subtype 10 (Disassoc) frame bursts. Reason codes and signal strength (RSSI) reported. \\
\bottomrule
\end{tabularx}
\caption{SOC Threat Classification and Verification Criteria.}
\end{table}

% =============================================================================
\section{Diagnostics and Troubleshooting}
% =============================================================================

\subsection{Common Issues and Solutions}

\paragraph{1. Serial Port Access Denied (\texttt{PermissionError} / \texttt{COM Port Busy}):}
\begin{itemize}[noitemsep]
    \item \textbf{Cause}: Arduino IDE Serial Monitor or another terminal instance has an exclusive lock on the COM port.
    \item \textbf{Solution}: Close the Serial Monitor window in Arduino IDE and restart \texttt{live\_collector.py}.
\end{itemize}

\paragraph{2. Wi-Fi Channel Mismatch Warning:}
\begin{itemize}[noitemsep]
    \item \textbf{Cause}: The Windows Hotspot adapter hopped to Channel 1 or 11, while the ESP32 defaulted to Channel 6.
    \item \textbf{Solution}: The updated collector automatically emits \texttt{\{"cmd":"set\_channel","channel":N\}} upon handshake, tuning the ESP32 in real time. If manual tuning is desired, set the default channel in line 42 of \texttt{esp32\_sensor\_node.ino} and reflash.
\end{itemize}

\paragraph{3. Memory Leak Prevention \& SQLite WAL Checkpointing:}
\begin{itemize}[noitemsep]
    \item The collector automatically runs \texttt{PRAGMA wal\_checkpoint(PASSIVE)} during periodic status heartbeats and \texttt{PRAGMA wal\_checkpoint(TRUNCATE)} on shutdown.
    \item This prevents virtual memory and Windows swap pagefile thrashing, maintaining RAM usage under 60 MB.
\end{itemize}

\paragraph{4. Terminating Background Zombie Processes:}
If an orphaned Streamlit or Python process locks port 8501 or the database, run the bundled PowerShell utility:
\begin{lstlisting}[language=bash, numbers=none]
powershell -ExecutionPolicy Bypass -File scripts/kill_zombies.ps1
\end{lstlisting}

% =============================================================================
\section{Conclusion}
% =============================================================================
This manual provides the foundational operational guidelines for the AI-Assisted IoT Wi-Fi Network Security System. With zero-footprint passive sniffing, autonomous database self-healing, and one-click frontend deployment, the system provides high-reliability, real-time intrusion monitoring for mission-critical IoT deployments.

\end{document}
